%=============================================================================
% INTRODUCTION
%=============================================================================
\section{Introduction}
\label{sec:intro}

\cA{Dengue is the fastest-spreading mosquito-borne viral disease globally, and Sri
Lanka faces recurring severe epidemics: the 2017 outbreak alone recorded over
186{,}000 cases and 440 deaths~\cite{tissera2020severe}. Roughly three weeks
separate infection from clinical diagnosis, so district-level forecasts issued
$H=3$ weeks ahead carry actionable lead time for pre-positioning clinical capacity
and intensifying vector control.}

\cB{Three properties of this task shape everything that follows. Transmission is
spatially coupled, which motivates a graph. The panel is small---25 districts over
459 weeks, against the hundreds of nodes typical in traffic benchmarks---so
parameters are cheap to overfit. And the target is severely heavy-tailed
(skewness 8.6; the top 1\% of district-weeks hold 18.3\% of all cases), which
makes squared error a measurement of a few dozen windows.}

\cC{A fourth property is easy to miss and dominates the rest: weekly incidence is
strongly autocorrelated. Cases at $t-1$ explain $r^2=0.85$ of cases at $t$. Any
model is therefore competing against a naive persistence forecast
$\hat{y}_{t+h}=y_t$ that is already close to optimal, and a benchmark that omits
that baseline cannot tell whether its architecture contributed anything. Weng
et al.~\cite{weng2024graph}, the reference benchmark for this dataset, omit it.}

\cD{We set out to add three contributions to an ST-GNN baseline---a learned
adaptive graph, a physics-informed loss derived from an SEIR--SEI compartmental
model, and GAN augmentation. Establishing what the baseline actually was turned
out to be the harder and more informative problem, and this paper reports that.
We reproduce the published results exactly, re-evaluate them honestly, identify a
data artifact that distorts the evaluation, and then test the mechanistic prior to
destruction---reporting not only that it fails but the quantity that makes failure
inevitable.}

\noindent\textbf{Contributions.}
\begin{itemize}\itemsep2pt
  \item \textbf{Exact reproduction, and what it reveals.} We reproduce five
        published ST-GNNs to within 7.3\% across twenty comparisons, two to four
        decimal places, and show that the reported column is computed with
        training data included and averaged per window rather than pooled. Naive
        persistence beats every model in it (\S\ref{sec:repro}).
  \item \textbf{An artifact-aware protocol.} We identify a 19$\times$ reporting
        backlog at week~395 spanning 18 of 25 districts. It falls in one fold's
        test split and no training split, contributes $\sim$90\% of that fold's
        squared error, and inverts fold difficulty. We score every arm and the
        baseline with and without it (\S\ref{sec:setup}).
  \item \textbf{A null result, stated as one.} Five verified architectures and
        four evidence-led increments, over 3 origins $\times$ 3 seeds. Against the
        artifact-free floor of 29.52, no increment differs from its own baseline
        at $p<0.05$ (\S\ref{sec:results}).
  \item \textbf{Why the mechanistic prior cannot help.} A renewal equation seeded
        by a validated SEIR--SEI generation interval reconstructs the record given
        oracle $R_t$ but loses as a forecaster. We trace this to a single
        measurement---$R_t$ is 26\% predictable, implying a 1.90$\times$ error per
        step---and rule out the two mechanisms that could have supplied the
        missing predictability (\S\ref{sec:physics}).
  \item \textbf{The target that is tractable.} Outbreaks are rankable three weeks
        ahead (AUC 0.826) even though they are not countable, and the mechanistic
        quantity is the largest single contributor to that ranking
        (\S\ref{sec:detection}).
\end{itemize}
